\section{映射}
\subsection{映射的概念}
\begin{definition}
%@see: 《Elements of Set Theory》 P42 Definition
设\(F\)是关系，如果\begin{equation*}
	(\forall x \in \dom F)(\exists! y)[\opair{x,y} \in F],
\end{equation*}
则称“关系\(F\)是一个\DefineConcept{映射}（function）”.
\end{definition}
可以从映射的定义中看出，虽然映射也是关系，
但映射有一般的关系所没有的特殊性质：
映射是\DefineConcept{单值的}（single-valued）.
换句话说，对于关系\(F\)，每个\(x\)可能对应若干个\(y\)；
但是，对于映射\(F\)，每个\(x\)就只对应一个\(y\).
我们可以把\(x\)与\(y\)这两个元素之间的对应关系
记为\(x \mapsto y\) \footnote{
	记号\(x \mapsto y\)
	读作“把\(x\)映成\(y\)的映射（the function that takes \(x\) to \(y\)）”.
}.

我们把使得\(xFy\)成立的\(y\)称为“\(x\)（在映射\(F\)下）的\DefineConcept{像}”
或“映射\(F\)在\(x\)的\DefineConcept{值}（the \emph{value} of \(F\) at \(x\)）”，
记为\(F(x)\)，即\begin{equation*}
	y = F(x);
\end{equation*}
称\(x\)为“\(y\)（在映射\(F\)下）的一个\DefineConcept{原像}”.
这里用的\(F(x)\)符号是欧拉提出的，
我们仅当\(F\)是一个映射且\(x\in\dom F\)时使用这个记号.
不过，我们也可以定义：\begin{equation*}
	F(x) \defeq \bigcup\Set{ y \given \opair{x,y} \in F }.
\end{equation*}
作为集合表达式，
右侧对任意\(F\)和\(x\)都有意义；
但本节涉及映射值的论断仍在\(F\)是映射且\(x\in\dom F\)的条件下使用\(F(x)\)，
不能用这个扩展约定代替一般关系的像或原像.

映射是如此重要，以至于各家对用于描述映射的术语没有达成统一.
以下是两种最常采用的术语.

设\(X,Y\)都是集合，
如果\(f\)是一个映射，且\(\dom f = X\)，\(\ran f \subseteq Y\)，
则称“\(f\)是从\(X\)到\(Y\)的\DefineConcept{映射}%
（\(f\) is a function \emph{from} \(X\) \emph{into} \(Y\)）”，
或称“\(f\)将\(X\)映射到\(Y\)里%
（\(f\) \emph{maps} \(X\) \emph{into} \(Y\)）”，
记作\begin{equation*}
	f\colon X \to Y.
\end{equation*}
如果还有\(\ran f = Y\)，
那么称“\(f\)是从\(X\)到\(Y\)上的映射%
（\(f\) is a function from \(X\) \emph{onto} \(Y\)）”，
或称“\(f\)将\(X\)映射到\(Y\)上%
（\(f\) \emph{maps} \(X\) \emph{onto} \(Y\)）”，
或称“\(f\)是\DefineConcept{满射}（surjective）”.
我们把集合\(Y\)称为“映射\(f\)的\DefineConcept{陪域}（codomain）”.
我们可以说“任意映射总将它的定义域映射到它的值域上”，
还可以说“任意映射总把它的定义域映射到以它的值域为子集的任意集合\(Y\)里”.
注意到两种说法的区别，“上”字和“里”字的选用，
不光取决于映射\(f\)本身，还取决于我们讨论的陪域\(Y\).

如果\begin{equation*}
	(\forall y \in \ran f)
	(\exists! x)
	[\opair{x,y} \in f],
\end{equation*}
那么称“映射\(f\)是\DefineConcept{一对一的}（one-to-one）”.

有时候我们希望把“一对一的”这个概念套用到一般的关系上，
它们往往不是映射，因此我们类比于“单值的”，创造出“单根的”这个概念.
\begin{definition}
%@see: 《Elements of Set Theory》 P43 Definition
如果集合\(R\)满足\begin{equation*}
	(\forall y \in \ran R)
	(\exists! x)
	[\opair{x,y} \in R],
\end{equation*}
则称“\(R\)是\DefineConcept{单根的}（single-rooted）”.
\end{definition}

因此，我们可以说，“一个映射是单根的”当且仅当“这个映射是一对一的”.

如果\begin{equation*}
	(\forall x_1, x_2 \in \dom f)
	[x_1 \neq x_2 \implies f(x_1) \neq f(x_2)],
\end{equation*}
那么称“\(f\)是\DefineConcept{单射}（injective）”.

由于映射本就是单值的，若它还是单根的，那么这个映射就是单射.
换句话说，“一对一的映射”和“单射”是相同的概念.

如果\(f\)既是单射，又是满射，
那么称“\(f\)是\DefineConcept{双射}（bijective）
或\DefineConcept{一一映射}”.

我们可以给出一个最平凡的一一映射.
\begin{definition}
设\(X\)是集合.
我们把\begin{equation*}
	\Set{ \opair{x,x} \given x \in X }
\end{equation*}称为
“（\(X\)上的）\DefineConcept{恒等映射}（the \emph{identity function} on \(X\)）
或\DefineConcept{恒同映射}”，
记作\(I_X\)或\(1_X\).
\end{definition}

\begin{definition}
设\(F\)是映射，
集合\(D \subseteq \dom F\)，
\(C \in \ran F\).
\begin{itemize}
	\item 如果\begin{equation*}
		(\forall x \in D)[F(x) = C],
	\end{equation*}
	则称“\(F\)在\(D\)上\DefineConcept{恒等于}~\(C\)”，
	记作\(F(x) \equiv C\ (x \in D)\)；
	当\(D=\dom F\)时，
	称“\(F\)~\DefineConcept{恒等于}~\(C\)”，
	记为\(F(x) \equiv C\).
	\item 如果\begin{equation*}
		(\exists x \in D)[F(x) \neq C],
	\end{equation*}
	则称“\(F\)在\(D\)上\DefineConcept{不恒等于}~\(C\)”，
	记作\(F(x) \not\equiv C\ (x \in D)\)；
	当\(D=\dom F\)时，
	称“\(F\)~\DefineConcept{不恒等于}~\(C\)”，
	记为\(F(x) \not\equiv C\).
	\item 如果\begin{equation*}
		(\forall x \in D)[F(x) \neq C],
	\end{equation*}
	则称“\(F\)在\(D\)上\DefineConcept{恒不等于}~\(C\)”，
	记作\(F(x) \neq C\ (x \in D)\)；
	当\(D=\dom F\)时，
	称“\(F\)~\DefineConcept{恒不等于}~\(C\)”，
	记为\(F(x) \neq C\).
\end{itemize}
\end{definition}

\begin{example}\label{example:映射.定义域中的元素在映射下的像一定属于值域}
    设\(G\)是映射.
    证明：
    \begin{equation*}
        x \in \dom G
        \implies
        G(x) \in \ran G.
    \end{equation*}
    \begin{proof}
        由\(x\in\dom G\)，
        存在唯一的\(y\)满足\(\opair{x,y}\in G\).
        这个\(y\)就是\(G(x)\).
        因此\(\opair{x,G(x)}\in G\)，
        由值域的定义可知\(G(x)\in\ran G\).
    \end{proof}
\end{example}

\begin{example}
    %@see: 《Elements of Set Theory》 P52 Exercise 11
    设\(F,G\)都是映射，
    \(\dom F = \dom G = X\)，且
    \begin{equation*}
        (\forall x \in X)[F(x) = G(x)].
    \end{equation*}
    证明：\(F=G\).
    \begin{proof}
        任取\(\opair{x,y}\in F\).
        由\(\dom F=X\)可知\(x\in X\)，
        并且\(y=F(x)=G(x)\)，
        所以\(\opair{x,y}\in G\).
        因此\(F\subseteq G\).
        交换\(F\)与\(G\)的地位，
        同理得到\(G\subseteq F\).
        两个映射都只含有序对，
        故由外延公理可得\(F=G\).
    \end{proof}
\end{example}

\begin{example}
    %@see: 《Elements of Set Theory》 P52 Exercise 12
    设\(f,g\)都是映射.
    证明：
    \begin{equation*}
        f \subseteq g
        \iff
        \dom f \subseteq \dom g
        \land
        (\forall x \in \dom f)
        [f(x) = g(x)].
    \end{equation*}
    \begin{proof}
        先证必要性.
        设\(f\subseteq g\)，
        任取\(x\in\dom f\).
        因为\(\opair{x,f(x)}\in f\subseteq g\)，
        所以\(x\in\dom g\).
        由\(g\)的单值性得\(g(x)=f(x)\).
        因此\(\dom f\subseteq\dom g\)，
        并且两个映射在\(\dom f\)上逐点相等.

        再证充分性.
        设右侧的两个条件成立，
        任取\(\opair{x,y}\in f\).
        于是\(x\in\dom f\subseteq\dom g\)，
        且\(y=f(x)=g(x)\)，
        故\(\opair{x,y}\in g\).
        这就证明了\(f\subseteq g\).
    \end{proof}
\end{example}

\begin{example}
    %@see: 《Elements of Set Theory》 P53 Exercise 13
    设\(f,g\)都是映射，\(f \subseteq g\)且\(\dom g \subseteq \dom f\).
    证明：\(f=g\).
    \begin{proof}
        由\(f\subseteq g\)及上一例可知
        \begin{align*}
            \dom f &\subseteq\dom g, \\
            f(x) &=g(x)\quad(x\in\dom f).
        \end{align*}
        结合假设\(\dom g\subseteq\dom f\)，
        得到\(\dom f=\dom g\).
        两个映射具有相同的定义域并且逐点相等，
        故\(f=g\).
    \end{proof}
\end{example}

\begin{example}
    %@see: 《Elements of Set Theory》 P53 Exercise 14
    设\(f,g\)都是映射.
    证明：
    \begin{enumerate}
        \item \(f \cap g\)是映射.
        \item \(f \cup g\)是映射的充分必要条件是
        \begin{equation*}
            (\forall x \in (\dom f)\cap(\dom g))[f(x)=g(x)].
        \end{equation*}
    \end{enumerate}
    \begin{proof}
        首先，
        \(f\cap g\)是关系.
        若\(\opair{x,y},\opair{x,z}\in f\cap g\)，
        则这两个有序对都属于\(f\)，
        由\(f\)的单值性可得\(y=z\).
        而对\(f\cap g\)的定义域中的每个元素，
        相应的像按定义存在，
        故\(f\cap g\)是映射.

        其次，
        假设\(f\cup g\)是映射.
        若\(x\in\dom f\cap\dom g\)，
        则\(\opair{x,f(x)}\)与\(\opair{x,g(x)}\)都属于\(f\cup g\).
        由单值性可得\(f(x)=g(x)\)，
        所以题给条件必要.

        反之，
        假设\(f\)与\(g\)在公共定义域上逐点相等.
        任取\(\opair{x,y},\opair{x,z}\in f\cup g\).
        若两者同属于\(f\)或同属于\(g\)，
        则由该映射的单值性得\(y=z\).
        否则交换\(y,z\)的名称后，
        可设\(\opair{x,y}\in f\)、\(\opair{x,z}\in g\).
        这时\(x\in\dom f\cap\dom g\)，
        从而\(y=f(x)=g(x)=z\).
        因此关系\(f\cup g\)是单值的，
        也就是映射.
    \end{proof}
\end{example}

\begin{example}
    %@see: 《Elements of Set Theory》 P53 Exercise 15
    设\(\mathscr{A}\)是一组映射，且
    \begin{equation*}
        (\forall f,g\in\mathscr{A})[f \subseteq g \lor g \subseteq f].
    \end{equation*}
    证明：\(\bigcup\mathscr{A}\)也是映射.
    \begin{proof}
        令\(H=\bigcup\mathscr{A}\).
        \(H\)的每个元素都属于某个映射，
        因而都是有序对，
        所以\(H\)是关系.
        若\(\opair{x,y},\opair{x,z}\in H\)，
        则存在\(f,g\in\mathscr{A}\)，
        使得\(\opair{x,y}\in f\)、\(\opair{x,z}\in g\).
        由假设，
        \(f\subseteq g\)或\(g\subseteq f\).
        因此这两个有序对同属于其中较大的那个映射，
        由其单值性可得\(y=z\).
        所以\(H\)是映射.
        当\(\mathscr{A}=\emptyset\)时，
        上述单值性条件真空成立，
        此时\(H=\emptyset\)是空映射.
    \end{proof}
\end{example}

\begin{example}
    %@see: 《Elements of Set Theory》 P53 Exercise 16
    证明：不存在一个集合，使得每个映射都属于它.
    \begin{proof}
        反设存在集合\(M\)包含每个映射.
        对任意集合\(a\)，
        由对集公理可以构造映射\(f_a=\Set{\opair{a,a}}\)，
        于是\(f_a\in M\).
        按有序对的定义，
        \(\opair{a,a}=\Set{\Set{a}}\)，
        故
        \begin{align*}
            \opair{a,a} &\in\bigcup M, \\
            \Set{a} &\in\bigcup\bigcup M, \\
            a &\in\bigcup\bigcup\bigcup M.
        \end{align*}
        令\(U=\bigcup\bigcup\bigcup M\)，
        则\(U\)是集合且包含每个集合.
        由子集公理构造\(D=\Set{x\in U\given x\notin x}\).
        因为\(D\)也是集合，
        所以\(D\in U\)，
        从而
        \begin{align*}
            D\in D &\iff D\in U\land D\notin D \\
            &\iff D\notin D.
        \end{align*}
        这与排中律和矛盾律相矛盾，
        所以这样的\(M\)不存在.
    \end{proof}
\end{example}

\begin{example}
    %@see: 《实变函数论（第三版）》（周民强） P14 思考题5.
    设\(f\colon X \to Y,
    g\colon Y \to X\).
    证明：若
    \begin{equation*}
        (\forall x \in X)[g(f(x)) = x],
    \end{equation*}
    则\(f\)是单射，\(g\)是满射.
    \begin{proof}
        若\(x_1,x_2\in X\)且\(f(x_1)=f(x_2)\)，
        则
        \begin{align*}
            x_1 &=g(f(x_1))=g(f(x_2))=x_2.
        \end{align*}
        故\(f\)是单射.
        另一方面，
        对任意\(x\in X\)，
        \(f(x)\in Y\)且\(g(f(x))=x\).
        所以\(x\in\ran g\)，
        进而\(X\subseteq\ran g\).
        结合\(g\colon Y\to X\)可得\(\ran g=X\)，
        故\(g\)是满射.
    \end{proof}
\end{example}

\subsection{逆，复合，限制，像，原像}
以下定义的操作通常用在映射上，有时候也用于关系，但也可以用于任意集合.
\begin{definition}\label{definition:映射.逆-复合-限制-像}
设\(A,F,G\)都是集合.
\begin{itemize}
	\item 称集合\begin{equation*}
		\Set*{ \opair{u,v} \given \opair{v,u} \in F }
	\end{equation*}为“\(F\)的\DefineConcept{逆}%
	（the \emph{inverse} of \(F\)）”，
	记作\(F^{-1}\).

	特别地，如果\(F\)和\(F^{-1}\)都是映射，
	则称“\(F^{-1}\)是\(F\)的\DefineConcept{逆映射}”.

	\item 称集合\begin{equation*}
		\Set*{ \opair{u,v} \given (\exists t)[\opair{u,t} \in G \land \opair{t,v} \in F] }
	\end{equation*}为“\(F\)和\(G\)的\DefineConcept{复合}%
	（the \emph{composition} of \(F\) and \(G\)）”，
	记作\(F \circ G\).

	\item 称集合\begin{equation*}
		\Set*{ \opair{u,v} \given \opair{u,v} \in F \land u \in A }
	\end{equation*}为“\(F\)在\(A\)上的\DefineConcept{限制}%
	（the \emph{restriction} of \(F\) to \(A\)）”，
	记作\(F \SetRestrict A\).
	%\cref{definition:子空间.拓扑子空间中的集族的限制}

	\item 称集合\begin{equation*}
		\Set*{ v \given (\exists u \in A)[\opair{u,v} \in F] }
	\end{equation*}为“\(A\)在\(F\)下的\DefineConcept{像}%
	（the \emph{image} of \(A\) \emph{under} \(F\)）”，
	记作\(F\ImageOfSetUnderRelation{A}\).
\end{itemize}
\end{definition}

当\(F\)是一个映射，且\(A \subseteq \dom F\)时，
\(F\ImageOfSetUnderRelation{A}\)这个概念可能更容易理解，
因为这时候\begin{equation*}
	F\ImageOfSetUnderRelation{A}
	= \Set{ F(u) \given u \in A }.
\end{equation*}

我们可以利用子集公理构造出上述定义下的所需集合的存在性.
特别地，\begin{gather*}
	F^{-1} \subseteq \ran F \times \dom F, \qquad
	F \circ G \subseteq \dom G \times \ran F, \\
	F \SetRestrict A \subseteq F, \qquad
	F\ImageOfSetUnderRelation{A} \subseteq \ran F.
\end{gather*}

例如，我们可以按如下方法正当化“关系\(F\)的逆”的定义：
根据子集公理，存在集合\(B\)，使得对于任意\(x\)，总有\begin{align*}
	x \in B
	&\iff
	[x \in \ran F \times \dom F]
	\land
	(\exists u)(\exists v)[x = \opair{u,v} \land \opair{v,u} \in F], \\
	&\iff
	(\exists u)(\exists v)[x = \opair{u,v} \land \opair{v,u} \in F].
\end{align*}
再根据外延公理，可以保证集合\(B\)的唯一性.
因此我们可以将集合\(B\)记为\(F^{-1}\).

\begin{definition}
%@see: 《点集拓扑讲义（第四版）》（熊金城） P22 定义1.5.4
设\(A,B,X\)都是集合，\(A \subset B\).
若映射\(F\colon A \to X\)和\(G\colon B \to X\)满足\(F \subset G\)，
则称“\(G\)是\(F\)在\(B\)上的\DefineConcept{扩张}”.
% 延拓
\end{definition}

\begin{theorem}
    \(F \SetRestrict \emptyset = \emptyset\).
    \begin{proof}
        按限制的定义，
        \(F\SetRestrict\emptyset\)中的有序对\(\opair{x,y}\)
        必须满足\(x\in\emptyset\).
        这样的\(x\)不存在，
        因此该限制没有元素，
        也就是\(F\SetRestrict\emptyset=\emptyset\).
    \end{proof}
\end{theorem}

\begin{theorem}
    \(F\ImageOfSetUnderRelation{A} = \ran(F \SetRestrict A)\).
    \begin{proof}
        任取\(y\)，
        由值域、限制和像的定义可得
        \begin{align*}
            y\in\ran(F\SetRestrict A)
            &\iff(\exists x)[\opair{x,y}\in F\SetRestrict A] \\
            &\iff(\exists x\in A)[\opair{x,y}\in F] \\
            &\iff y\in F\ImageOfSetUnderRelation{A}.
        \end{align*}
        由外延公理得到所需等式.
    \end{proof}
\end{theorem}

\begin{example}
    %@see: 《Elements of Set Theory》 P53 Exercise 18.
    设
    \begin{equation*}
        R = \Set{
        \opair{0,1},
        \opair{0,2},
        \opair{0,3},
        \opair{1,2},
        \opair{1,3},
        \opair{2,3}
        }.
    \end{equation*}
    计算\(R \circ R,
    R \SetRestrict \{1\},
    R^{-1} \SetRestrict \{1\},
    R\ImageOfSetUnderRelation{\{1\}},
    R^{-1}\ImageOfSetUnderRelation{\{1\}}\).
    \begin{solution}
        在复合\(R\circ R\)中，
        中间元素可以是\(1\)或\(2\).
        中间元素为\(1\)时得到\(\opair{0,2},\opair{0,3}\)；
        中间元素为\(2\)时得到\(\opair{0,3},\opair{1,3}\).
        去除重复的有序对，
        并分别按首坐标取限制、按末坐标取像，
        得到
        \begin{align*}
            R\circ R &=\Set{\opair{0,2},\opair{0,3},\opair{1,3}}, \\
            R\SetRestrict\Set{1} &=\Set{\opair{1,2},\opair{1,3}}, \\
            R^{-1}\SetRestrict\Set{1} &=\Set{\opair{1,0}}, \\
            R\ImageOfSetUnderRelation{\Set{1}} &=\Set{2,3}, \\
            R^{-1}\ImageOfSetUnderRelation{\Set{1}} &=\Set{0}.
        \end{align*}
        % 精确有限集合核算见随本次改动提供的 Python 与 Mathematica 验证脚本.
    \end{solution}
\end{example}

\begin{example}
    %@see: 《Elements of Set Theory》 P53 Exercise 19.
    设
    \begin{equation*}
        A = \Set{
        \opair{\emptyset,\{\emptyset,\{\emptyset\}\}},
        \opair{\{\emptyset\},\emptyset}
        }.
    \end{equation*}
    计算\(A(\emptyset),
    A\ImageOfSetUnderRelation{\emptyset},
    A\ImageOfSetUnderRelation{\{\emptyset\}},
    A\ImageOfSetUnderRelation{\{\emptyset,\{\emptyset\}\}},
    A^{-1},
    A \circ A,
    A \SetRestrict \emptyset,
    A \SetRestrict \{\emptyset\},
    A \SetRestrict \{\emptyset,\{\emptyset\}\},
    \bigcup\bigcup A\).
    \begin{solution}
        为清楚区分元素与单元素集，
        暂记\(e=\emptyset\)、\(s=\Set{\emptyset}\)、\(t=\Set{\emptyset,\Set{\emptyset}}\).
        这三个集合两两不同，
        而题给映射就是\(A=\Set{\opair{e,t},\opair{s,e}}\).
        因此\(\dom A=\Set{e,s}\)，
        且\(A(e)=t\)、\(A(s)=e\).
        由于\(t\notin\dom A\)，
        复合\(A\circ A\)仅在\(s\)处有定义，
        其值为\(A(A(s))=A(e)=t\).
        逐项计算得到
        \begin{align*}
            A(\emptyset) &=t, \\
            A\ImageOfSetUnderRelation{\emptyset} &=\emptyset, \\
            A\ImageOfSetUnderRelation{\Set{\emptyset}} &=\Set{t}, \\
            A\ImageOfSetUnderRelation{\Set{\emptyset,\Set{\emptyset}}} &=\Set{t,e}, \\
            A^{-1} &=\Set{\opair{t,e},\opair{e,s}}, \\
            A\circ A &=\Set{\opair{s,t}}, \\
            A\SetRestrict\emptyset &=\emptyset, \\
            A\SetRestrict\Set{\emptyset} &=\Set{\opair{e,t}}, \\
            A\SetRestrict\Set{\emptyset,\Set{\emptyset}} &=A.
        \end{align*}
        最后，
        按\(\opair{u,v}=\Set{\Set{u},\Set{u,v}}\)展开，
        可得
        \begin{align*}
            \bigcup A &=\Set{\Set{e},\Set{e,t},\Set{s},\Set{s,e}}, \\
            \bigcup\bigcup A &=\Set{e,s,t} \\
            &=\Set{\emptyset,\Set{\emptyset},\Set{\emptyset,\Set{\emptyset}}}.
        \end{align*}
        特别地，
        \(A\ImageOfSetUnderRelation{\emptyset}\)没有元素，
        而\(A\ImageOfSetUnderRelation{\Set{\emptyset}}=\Set{A(\emptyset)}\)是单元素集，
        二者不能混淆.
    \end{solution}
\end{example}

\begin{example}
    %@see: 《Elements of Set Theory》 P53 Exercise 20.
    证明：\(F \SetRestrict A = F \cap (A \times \ran F)\).
    \begin{proof}
        等式两边都是只含有序对的关系.
        对任意\(x,y\)，
        若\(\opair{x,y}\in F\)，
        则\(y\in\ran F\).
        因此
        \begin{align*}
            \opair{x,y}\in F\SetRestrict A
            &\iff\opair{x,y}\in F\land x\in A \\
            &\iff\opair{x,y}\in F\land x\in A\land y\in\ran F \\
            &\iff\opair{x,y}\in F\cap(A\times\ran F).
        \end{align*}
        由外延公理可得结论，
        此处不要求\(F\)本身是映射.
    \end{proof}
\end{example}

\begin{definition}
    设\(F\)是关系，
    \(A\)是集合.
    称集合
    \begin{align*}
        F^{-1}\ImageOfSetUnderRelation{A}
        &\defeq\Set{x\in\dom F\given(\exists y\in A)[\opair{x,y}\in F]},
    \end{align*}
    为“集合\(A\)在关系\(F\)下的\DefineConcept{原像}%
    （the \emph{inverse image} of \(A\) under \(F\)）”.
    这正是\(A\)在逆关系\(F^{-1}\)下的像.
    当\(F\)是映射时，
    单值性才使上述定义化为
    \begin{align*}
        F^{-1}\ImageOfSetUnderRelation{A}
        &=\Set{x\in\dom F\given F(x)\in A}.
    \end{align*}
\end{definition}

一般来说，
一个映射的逆不一定是映射.
例如，
取
\begin{align*}
    F &=\Set{\opair{1,1},\opair{2,1}}, \\
    F^{-1} &=\Set{\opair{1,1},\opair{1,2}}.
\end{align*}
\(F\)是映射，
但\(F^{-1}\)将同一个元素\(1\)对应到不同的元素\(1,2\)，
所以\(F^{-1}\)不是映射.

\begin{theorem}\label{theorem:集合论.关系的逆的定义域值域以及关系的二重逆}
    %@see: 《Elements of Set Theory》 P46 Theorem 3E
    设\(F\)是集合，则有
    \begin{gather}
        \dom F^{-1} = \ran F, \\
        \ran F^{-1} = \dom F.
    \end{gather}

    如果\(F\)是关系，则有
    \begin{equation}
        (F^{-1})^{-1} = F.
    \end{equation}
    \begin{proof}
        对任意\(y\)与\(x\)，
        分别有
        \begin{align*}
            y\in\dom F^{-1}
            &\iff(\exists u)[\opair{y,u}\in F^{-1}] \\
            &\iff(\exists u)[\opair{u,y}\in F] \\
            &\iff y\in\ran F, \\
            x\in\ran F^{-1}
            &\iff(\exists v)[\opair{v,x}\in F^{-1}] \\
            &\iff(\exists v)[\opair{x,v}\in F] \\
            &\iff x\in\dom F.
        \end{align*}
        由外延公理得到前两个等式.

        若\(F\)是关系，
        则\(F\)与\((F^{-1})^{-1}\)都只含有序对.
        对任意\(x,y\)，
        有
        \begin{align*}
            \opair{x,y}\in(F^{-1})^{-1}
            &\iff\opair{y,x}\in F^{-1} \\
            &\iff\opair{x,y}\in F.
        \end{align*}
        故\((F^{-1})^{-1}=F\).
        对一般集合，
        二重取逆只保留其中的有序对，
        因此最后这个等式不能无条件推广.
    \end{proof}
\end{theorem}

\begin{theorem}\label{theorem:集合论.关系及其逆是映射的充分必要条件}
    %@see: 《Elements of Set Theory》 P46 Theorem 3F
    设\(F\)是集合，则“\(F^{-1}\)是映射”的充分必要条件是：\(F\)是单根的.

    设\(F\)是关系，则“\(F\)是映射”的充分必要条件是：\(F^{-1}\)是单根的.
    \begin{proof}
        对任意集合\(F\)，
        其逆\(F^{-1}\)总是关系，
        并且\(\dom F^{-1}=\ran F\).
        由定义可得
        \begin{align*}
            F^{-1}\text{是映射}
            &\iff(\forall y\in\ran F)(\exists!x)[\opair{y,x}\in F^{-1}] \\
            &\iff(\forall y\in\ran F)(\exists!x)[\opair{x,y}\in F] \\
            &\iff F\text{是单根的}.
        \end{align*}

        若\(F\)还是关系，
        则\(F=(F^{-1})^{-1}\).
        在刚证明的等价关系中以\(F^{-1}\)代替\(F\)，
        得到
        \begin{align*}
            F\text{是映射}
            &\iff(F^{-1})^{-1}\text{是映射} \\
            &\iff F^{-1}\text{是单根的}.
        \end{align*}
    \end{proof}
\end{theorem}

\begin{theorem}\label{theorem:集合论.逆映射的计算}
    %@see: 《Elements of Set Theory》 P46 Theorem 3G
    设\(F\)是单射.
    \begin{enumerate}
        \item 如果\(x \in \dom F\)，那么
        \begin{equation*}
            F^{-1}(F(x)) = x.
        \end{equation*}

        \item 如果\(y \in \ran F\)，那么
        \begin{equation*}
            F(F^{-1}(y)) = y.
        \end{equation*}
    \end{enumerate}
    \begin{proof}
        由于\(F\)是单射，
        它既是单值的又是单根的，
        所以\(F^{-1}\)也是映射，
        且\(\dom F^{-1}=\ran F\).

        若\(x\in\dom F\)，
        则\(\opair{x,F(x)}\in F\)，
        从而\(\opair{F(x),x}\in F^{-1}\).
        由\(F^{-1}\)的单值性可得\(F^{-1}(F(x))=x\).

        若\(y\in\ran F\)，
        则\(\opair{y,F^{-1}(y)}\in F^{-1}\)，
        从而\(\opair{F^{-1}(y),y}\in F\).
        由\(F\)的单值性可得\(F(F^{-1}(y))=y\).
        这也同时验证了两式中每次求值均在相应的定义域内.
    \end{proof}
\end{theorem}

\begin{theorem}\label{theorem:集合论.映射的复合也是映射}
    %@see: 《Elements of Set Theory》 P47 Theorem 3H
    设\(F,G\)都是映射，则\(F \circ G\)是映射，且
    \begin{gather*}
        \dom(F \circ G)
        = \Set*{ x \in \dom G \given G(x) \in \dom F }, \\
        (\forall x \in \dom(F \circ G))
        [(F \circ G)(x) = F(G(x))].
    \end{gather*}
    \begin{proof}
        \(F\circ G\)按定义是关系.
        若\(\opair{x,y},\opair{x,z}\in F\circ G\)，
        则存在\(p,q\)满足
        \begin{align*}
            \opair{x,p}\in G &\land\opair{p,y}\in F, \\
            \opair{x,q}\in G &\land\opair{q,z}\in F.
        \end{align*}
        由\(G\)的单值性可得\(p=q\)，
        再由\(F\)的单值性可得\(y=z\).
        因此\(F\circ G\)是映射.

        若\(x\in\dom(F\circ G)\)，
        则存在\(t,y\)使得\(\opair{x,t}\in G\)且\(\opair{t,y}\in F\).
        故\(x\in\dom G\)，
        且\(G(x)=t\in\dom F\).
        反之，
        若\(x\in\dom G\)且\(G(x)\in\dom F\)，
        则
        \begin{align*}
            \opair{x,G(x)} &\in G, \\
            \opair{G(x),F(G(x))} &\in F.
        \end{align*}
        于是\(\opair{x,F(G(x))}\in F\circ G\).
        这既证明了所述定义域等式，
        也证明了\((F\circ G)(x)=F(G(x))\).
    \end{proof}
\end{theorem}

映射的复合有顺序之分.
按照本节的关系运算定义，
\(f\circ g\)与\(g\circ f\)总能构成关系，
并且当\(f,g\)是映射时两者也都是映射，
但它们的定义域可能缩小甚至为空.
若要求\(f\circ g\)定义在整个\(\dom g\)上，
还需要\(\ran g\subseteq\dom f\)；
反向复合则需要检查\(\ran f\subseteq\dom g\).
即便两个条件都满足，
两种复合也未必相同.

\begin{example}\label{example:映射.映射的复合适合结合律}
    %@see: 《Elements of Set Theory》 P53 Exercise 21.
    设\(R,S,T\)都是集合.
    证明：\((R \circ S) \circ T = R \circ (S \circ T)\).
    \begin{proof}
        两边都是关系，
        对任意有序对\(\opair{x,y}\)，
        有
        \begin{align*}
            \opair{x,y}\in(R\circ S)\circ T
            &\iff(\exists u)[\opair{x,u}\in T\land\opair{u,y}\in R\circ S] \\
            &\iff(\exists u)(\exists v)
            [\opair{x,u}\in T\land\opair{u,v}\in S\land\opair{v,y}\in R] \\
            &\iff(\exists v)[\opair{x,v}\in S\circ T\land\opair{v,y}\in R] \\
            &\iff\opair{x,y}\in R\circ(S\circ T).
        \end{align*}
        由外延公理得到结合律.
    \end{proof}
\end{example}

\begin{example}\label{example:映射.逆映射与映射的复合}
    %@see: 《Elements of Set Theory》 P47 Example
    设\(G\)是单射，
    \(I_{\dom G}\)是\(\dom G\)上的恒等映射.
    证明：\(G^{-1} \circ G = I_{\dom G}\).
    \begin{proof}
        由于\(G\)是单射，
        \(G^{-1}\)是映射，
        其定义域为\(\ran G\).
        对每个\(x\in\dom G\)，
        都有\(G(x)\in\ran G=\dom G^{-1}\).
        由\cref{theorem:集合论.映射的复合也是映射}可得
        \begin{align*}
            \dom(G^{-1}\circ G)
            &=\Set{x\in\dom G\given G(x)\in\dom G^{-1}} \\
            &=\dom G.
        \end{align*}
        再由\cref{theorem:集合论.逆映射的计算}，
        对每个\(x\in\dom G\)有
        \begin{align*}
            (G^{-1}\circ G)(x) &=G^{-1}(G(x))=x=I_{\dom G}(x).
        \end{align*}
        所以两个映射具有相同的定义域和相同的值，
        故\(G^{-1}\circ G=I_{\dom G}\).
    \end{proof}
\end{example}
\begin{example}\label{example:映射.映射与逆映射的复合}
    %@see: 《Elements of Set Theory》 P53 Exercise 25.
    设\(G\)是映射，
    \(I_{\ran G}\)是\(\ran G\)上的恒等映射.
    证明：\(G \circ G^{-1} = I_{\ran G}\).
    \begin{proof}
        这里不要求\(G^{-1}\)是映射，
        只需使用关系的复合.
        对任意\(u,v\)，
        有
        \begin{align*}
            \opair{u,v}\in G\circ G^{-1}
            &\iff(\exists t)[\opair{u,t}\in G^{-1}\land\opair{t,v}\in G] \\
            &\iff(\exists t)[\opair{t,u}\in G\land\opair{t,v}\in G] \\
            &\iff u=v\land u\in\ran G \\
            &\iff\opair{u,v}\in I_{\ran G}.
        \end{align*}
        第三个等价的正向由\(G\)的单值性得到；
        反向则由\(u\in\ran G\)选取满足\(\opair{t,u}\in G\)的一个\(t\)，
        再利用\(u=v\)得到.
        两边都是关系，
        故由外延公理得到结论.
    \end{proof}
\end{example}
\begin{example}\label{example:映射.可逆映射是双射}
    %@see: 《高等代数学（第四版）》（谢启鸿 姚慕生 吴泉水） P184 命题4.1.1
    设映射\(f\colon A \to B\)，
    \(I_A\)是\(A\)上的恒等映射，
    \(I_B\)是\(B\)上的恒等映射.
    证明：如果存在映射\(g\colon B \to A\)满足
    \begin{equation*}
        g \circ f = I_A,
        \qquad
        f \circ g = I_B,
    \end{equation*}
    则\(f\)是双射，且\(g = f^{-1}\).
    \begin{proof}
        若\(a_1,a_2\in A\)且\(f(a_1)=f(a_2)\)，
        则
        \begin{align*}
            a_1 &=g(f(a_1))=g(f(a_2))=a_2.
        \end{align*}
        因此\(f\)是单射.
        对任意\(b\in B\)，
        \(g(b)\in A\)且\(f(g(b))=b\)，
        所以\(b\in\ran f\).
        结合\(\ran f\subseteq B\)，
        可知\(f\)是满射，
        因而是双射.

        最后证明\(g=f^{-1}\).
        任取\(b\in B\)，
        由\(f(g(b))=b\)可得\(\opair{g(b),b}\in f\)，
        于是\(\opair{b,g(b)}\in f^{-1}\).
        由于\(f\)是单射，
        \(f^{-1}\)是定义在\(B\)上的映射，
        故其单值性给出\(g(b)=f^{-1}(b)\).
        两个映射的定义域均为\(B\)且逐点相等，
        因此\(g=f^{-1}\).
    \end{proof}
\end{example}

\begin{theorem}
    %@see: 《Elements of Set Theory》 P65 Exercise 53.
    设\(R,S\)是集合，那么
    \begin{gather}
        (R \cup S)^{-1} = R^{-1} \cup S^{-1},
        \label{equation:集合论.并的逆等于逆的并} \\
        (R \cap S)^{-1} = R^{-1} \cap S^{-1},
        \label{equation:集合论.交的逆等于逆的交} \\
        (R - S)^{-1} = R^{-1} - S^{-1}.
        \label{equation:集合论.差的逆等于逆的差}
    \end{gather}
    \begin{proof}
        每个等式的两边都只含有序对.
        对任意\(x,y\)，
        由取逆的定义分别得到
        \begin{align*}
            \opair{x,y}\in(R\cup S)^{-1}
            &\iff\opair{y,x}\in R\lor\opair{y,x}\in S \\
            &\iff\opair{x,y}\in R^{-1}\cup S^{-1}, \\
            \opair{x,y}\in(R\cap S)^{-1}
            &\iff\opair{y,x}\in R\land\opair{y,x}\in S \\
            &\iff\opair{x,y}\in R^{-1}\cap S^{-1}, \\
            \opair{x,y}\in(R-S)^{-1}
            &\iff\opair{y,x}\in R\land\opair{y,x}\notin S \\
            &\iff\opair{x,y}\in R^{-1}-S^{-1}.
        \end{align*}
        由外延公理可得三个等式.
    \end{proof}
\end{theorem}

\begin{theorem}\label{theorem:集合论.复合的逆}
    %@see: 《Elements of Set Theory》 P47 Theorem 3I
    设\(F,G\)都是集合，那么
    \begin{equation*}
        (F \circ G)^{-1} = G^{-1} \circ F^{-1}.
    \end{equation*}
    \begin{proof}
        两边都是关系，
        并且对任意\(x,y\)有
        \begin{align*}
            \opair{x,y}\in(F\circ G)^{-1}
            &\iff\opair{y,x}\in F\circ G \\
            &\iff(\exists t)[\opair{y,t}\in G\land\opair{t,x}\in F] \\
            &\iff(\exists t)[\opair{x,t}\in F^{-1}\land\opair{t,y}\in G^{-1}] \\
            &\iff\opair{x,y}\in G^{-1}\circ F^{-1}.
        \end{align*}
        由外延公理得到结论.
    \end{proof}
\end{theorem}

下面的结论对任意集合都成立，
可以直接从复合与取逆的定义证明.
\begin{proposition}\label{theorem:集合论.复合的逆.推论1}
    设\(F\)是集合，那么
    \begin{equation*}
        (F^{-1} \circ F)^{-1} = F^{-1} \circ F.
    \end{equation*}
    \begin{proof}
        对任意集合\(F\)，
        \(F^{-1}\circ F\)都是关系.
        由定义可得
        \begin{align*}
            \opair{x,y}\in(F^{-1}\circ F)^{-1}
            &\iff\opair{y,x}\in F^{-1}\circ F \\
            &\iff(\exists t)[\opair{y,t}\in F\land\opair{x,t}\in F] \\
            &\iff(\exists t)[\opair{x,t}\in F\land\opair{y,t}\in F] \\
            &\iff\opair{x,y}\in F^{-1}\circ F.
        \end{align*}
        由外延公理可得等式.
        这里直接使用有序对，
        没有对一般集合\(F\)误用\((F^{-1})^{-1}=F\).
    \end{proof}
\end{proposition}

\begin{axiom}[选择公理(第一种形式)]
%@see: 《Elements of Set Theory》 P49 Axiom of Choice (first form)
对于任意关系\(R\)，存在映射\(H\)，满足\begin{equation*}
	H \subseteq R,
	\quad\text{且}\quad
	\dom H = \dom R.
\end{equation*}
\end{axiom}

\begin{theorem}
    %@see: 《Elements of Set Theory》 P48 Theorem 3J
    设映射\(F\colon A \to B\)，其中\(A\)是非空集合.
    \begin{enumerate}
        \item “存在映射\(G\colon B \to A\)（称其为\DefineConcept{左逆}），
        使得\(G \circ F\)是\(A\)上的恒等映射\(I_A\)”是“\(F\)是单射”的充分必要条件.

        \item “存在映射\(H\colon B \to A\)（称其为\DefineConcept{右逆}），
        使得\(F \circ H\)是\(B\)上的恒等映射\(I_B\)”是“\(F\)是满射”的充分必要条件.
    \end{enumerate}
    \begin{proof}
        首先证明关于左逆的断言.
        若存在\(G\colon B\to A\)满足\(G\circ F=I_A\)，
        则对满足\(F(x)=F(y)\)的\(x,y\in A\)有
        \begin{align*}
            x &=G(F(x))=G(F(y))=y.
        \end{align*}
        所以\(F\)是单射.
        反之，
        若\(F\)是单射，
        则\(F^{-1}\colon\ran F\to A\)是映射.
        利用\(A\neq\emptyset\)取定\(a_0\in A\)，
        构造
        \begin{align*}
            G &=F^{-1}\cup\bigl((B-\ran F)\times\Set{a_0}\bigr).
        \end{align*}
        右边的两个映射具有不相交的定义域，
        其并是定义在\(B\)上的映射，
        且值均属于\(A\).
        对每个\(x\in A\)有
        \begin{align*}
            G(F(x)) &=F^{-1}(F(x))=x.
        \end{align*}
        因此\(G\circ F=I_A\).
        这个构造只选取一个固定元素，
        不需要选择公理.

        其次证明关于右逆的断言.
        若存在\(H\colon B\to A\)满足\(F\circ H=I_B\)，
        则对任意\(y\in B\)，
        \(H(y)\in A\)且\(F(H(y))=y\).
        所以\(B\subseteq\ran F\)，
        结合\(\ran F\subseteq B\)可知\(F\)是满射.
        反之，
        设\(\ran F=B\).
        对关系\(F^{-1}\)应用选择公理，
        得到映射\(H\)满足
        \begin{align*}
            H &\subseteq F^{-1}, \\
            \dom H &=\dom F^{-1}=B.
        \end{align*}
        又有\(\ran H\subseteq\ran F^{-1}=A\)，
        故\(H\colon B\to A\).
        对任意\(y\in B\)，
        \(\opair{y,H(y)}\in H\subseteq F^{-1}\)，
        所以\(\opair{H(y),y}\in F\)，
        进而\(F(H(y))=y\).
        因此\(F\circ H=I_B\).
    \end{proof}
\end{theorem}

\begin{theorem}
    %@see: 《Elements of Set Theory》 P48 Theorem 3K
    设\(A,B,F\)都是集合.
    \begin{enumerate}
        \item 并的像是像的并：
        \begin{gather}
            F\ImageOfSetUnderRelation{A \cup B}
            = F\ImageOfSetUnderRelation{A} \cup F\ImageOfSetUnderRelation{B},
            \label{equation:集合论.并的像与像的并的关系1} \\
            F\ImageOfSetUnderRelation{\bigcup A}
            = \bigcup\Set{ F\ImageOfSetUnderRelation{a} \given a \in A }.
            \label{equation:集合论.并的像与像的并的关系2}
        \end{gather}

        \item 交的像包含于像的交：
        \begin{gather}
            F\ImageOfSetUnderRelation{A \cap B}
            \subseteq F\ImageOfSetUnderRelation{A} \cap F\ImageOfSetUnderRelation{B},
            \label{equation:集合论.交的像与像的交的关系1} \\
            F\ImageOfSetUnderRelation{\bigcap A}
            \subseteq \bigcap\Set{ F\ImageOfSetUnderRelation{a} \given a \in A }.
            \label{equation:集合论.交的像与像的交的关系2}
            \quad(A \neq \emptyset)
        \end{gather}
        若\(F\)是单根的，则以上两式取“=”号.

        \item 差的像包含像的差：
        \begin{equation}
            F\ImageOfSetUnderRelation{A} - F\ImageOfSetUnderRelation{B}
            \subseteq F\ImageOfSetUnderRelation{A-B}.
            \label{equation:集合论.差的像与像的差的关系}
        \end{equation}
        若\(F\)是单根的，则上式取“=”号.
    \end{enumerate}
    \begin{proof}
        所出现的像族都可以从\(\Powerset(\ran F)\)中用子集公理分离得到，
        因而确实是集合.
        下面分别验证各式.

        对于两个集合的并，
        任取\(y\)，
        有
        \begin{align*}
            y\in F\ImageOfSetUnderRelation{A\cup B}
            &\iff(\exists x\in A\cup B)[\opair{x,y}\in F] \\
            &\iff(\exists x\in A)[\opair{x,y}\in F]
            \lor(\exists x\in B)[\opair{x,y}\in F] \\
            &\iff y\in F\ImageOfSetUnderRelation{A}\cup F\ImageOfSetUnderRelation{B}.
        \end{align*}
        对于任意集合族的并，
        同样有
        \begin{align*}
            y\in F\ImageOfSetUnderRelation{\bigcup A}
            &\iff(\exists x)[x\in\bigcup A\land\opair{x,y}\in F] \\
            &\iff(\exists a\in A)(\exists x\in a)[\opair{x,y}\in F] \\
            &\iff y\in\bigcup\Set{F\ImageOfSetUnderRelation{a}\given a\in A}.
        \end{align*}
        由外延公理得到两个并的等式.
        后一式在\(A=\emptyset\)时也成立，
        此时两边均为空集.

        对于两个集合的交，
        若\(y\in F\ImageOfSetUnderRelation{A\cap B}\)，
        则存在\(x\in A\cap B\)满足\(\opair{x,y}\in F\).
        同一个\(x\)分别见证\(y\in F\ImageOfSetUnderRelation{A}\)
        和\(y\in F\ImageOfSetUnderRelation{B}\)，
        故所述包含关系成立.
        若\(F\)是单根的，
        反过来设\(y\)属于这两个像的交，
        则存在\(u\in A\)、\(v\in B\)使得\(\opair{u,y},\opair{v,y}\in F\).
        单根性给出\(u=v\in A\cap B\)，
        故\(y\in F\ImageOfSetUnderRelation{A\cap B}\)，
        从而得到等式.

        现在设\(A\neq\emptyset\).
        若\(y\in F\ImageOfSetUnderRelation{\bigcap A}\)，
        取\(x\in\bigcap A\)使得\(\opair{x,y}\in F\).
        对每个\(a\in A\)，
        都有\(x\in a\)，
        从而\(y\in F\ImageOfSetUnderRelation{a}\).
        于是\(y\in\bigcap\Set{F\ImageOfSetUnderRelation{a}\given a\in A}\).
        若\(F\)单根，
        反过来设\(y\)属于右侧的交.
        取定一个\(a_0\in A\)，
        再取\(x_0\in a_0\)满足\(\opair{x_0,y}\in F\).
        对任意\(a\in A\)，
        由\(y\in F\ImageOfSetUnderRelation{a}\)存在\(x\in a\)满足\(\opair{x,y}\in F\).
        单根性保证\(x=x_0\)，
        所以\(x_0\in a\).
        因此\(x_0\in\bigcap A\)，
        从而\(y\in F\ImageOfSetUnderRelation{\bigcap A}\).
        这里使用一个固定的原像\(x_0\)，
        不需要对各个\(a\)同时作选择.

        最后，
        若\(y\in F\ImageOfSetUnderRelation{A}-F\ImageOfSetUnderRelation{B}\)，
        则存在\(x\in A\)满足\(\opair{x,y}\in F\).
        若\(x\in B\)，
        就会推出\(y\in F\ImageOfSetUnderRelation{B}\)，
        与假设矛盾.
        故\(x\in A-B\)，
        于是\(y\in F\ImageOfSetUnderRelation{A-B}\).
        若\(F\)单根，
        反过来设\(y\in F\ImageOfSetUnderRelation{A-B}\)，
        取\(x\in A-B\)满足\(\opair{x,y}\in F\).
        这已经给出\(y\in F\ImageOfSetUnderRelation{A}\).
        若还存在\(u\in B\)满足\(\opair{u,y}\in F\)，
        则单根性给出\(u=x\)，
        与\(x\notin B\)矛盾.
        故\(y\notin F\ImageOfSetUnderRelation{B}\)，
        从而差的包含关系也成为等式.
    \end{proof}
\end{theorem}

\begin{corollary}
    %@see: 《Elements of Set Theory》 P48 Corollary 3L
    设\(G\)是映射，\(A,B\)都是集合.
    \begin{gather}
        G^{-1}\ImageOfSetUnderRelation{\bigcup A} = \bigcup\Set*{ G^{-1}\ImageOfSetUnderRelation{a} \given a \in A },
        \label{equation:集合论.并的原像与原像的并的关系} \\
        G^{-1}\ImageOfSetUnderRelation{\bigcap A} = \bigcap\Set*{ G^{-1}\ImageOfSetUnderRelation{a} \given a \in A }, \quad A \neq \emptyset,
        \label{equation:集合论.交的原像与原像的交的关系} \\
        G^{-1}\ImageOfSetUnderRelation{A - B} = G^{-1}\ImageOfSetUnderRelation{A} - G^{-1}\ImageOfSetUnderRelation{B}.
        \label{equation:集合论.差的原像与原像的差的关系}
    \end{gather}
    \begin{proof}
        由\cref{theorem:集合论.关系及其逆是映射的充分必要条件}，
        映射\(G\)的逆关系\(G^{-1}\)是单根的.
        把上一条定理中的\(F\)替换为\(G^{-1}\)，
        并的像的等式直接给出\cref{equation:集合论.并的原像与原像的并的关系}；
        当\(A\neq\emptyset\)时，
        交的像的等式给出\cref{equation:集合论.交的原像与原像的交的关系}；
        差的像的等式给出\cref{equation:集合论.差的原像与原像的差的关系}.
        这里的\(G^{-1}\ImageOfSetUnderRelation{C}\)表示逆关系作用于集合\(C\)的像，
        不要求\(G^{-1}\)是逆映射，
        也不要求\(G\)是单射或满射.
    \end{proof}
\end{corollary}

\begin{example}
    %@see: 《Elements of Set Theory》 P53 Exercise 22.(a)
    证明：
    \begin{equation}
        A \subseteq B \implies F\ImageOfSetUnderRelation{A} \subseteq F\ImageOfSetUnderRelation{B}.
    \end{equation}
    \begin{proof}
        任取\(y\in F\ImageOfSetUnderRelation{A}\)，
        则存在\(x\in A\)满足\(\opair{x,y}\in F\).
        由\(A\subseteq B\)可知\(x\in B\)，
        所以同一个\(x\)见证\(y\in F\ImageOfSetUnderRelation{B}\).
        因此\(F\ImageOfSetUnderRelation{A}\subseteq F\ImageOfSetUnderRelation{B}\).
    \end{proof}
\end{example}

\begin{example}
    %@see: 《Elements of Set Theory》 P53 Exercise 22.(b)
    证明：
    \begin{equation}
        (F \circ G)\ImageOfSetUnderRelation{A}
        = F\ImageOfSetUnderRelation{G\ImageOfSetUnderRelation{A}}.
    \end{equation}
    \begin{proof}
        任取\(y\)，
        由像与复合的定义可得
        \begin{align*}
            y\in(F\circ G)\ImageOfSetUnderRelation{A}
            &\iff(\exists x\in A)[\opair{x,y}\in F\circ G] \\
            &\iff(\exists x\in A)(\exists t)
            [\opair{x,t}\in G\land\opair{t,y}\in F] \\
            &\iff(\exists t\in G\ImageOfSetUnderRelation{A})[\opair{t,y}\in F] \\
            &\iff y\in F\ImageOfSetUnderRelation{G\ImageOfSetUnderRelation{A}}.
        \end{align*}
        由外延公理得到等式.
    \end{proof}
\end{example}

\begin{example}
    %@see: 《Elements of Set Theory》 P53 Exercise 22.(c)
    证明：
    \begin{equation}
        Q \SetRestrict (A \cup B)
        = (Q \SetRestrict A)\cup(Q \SetRestrict B).
    \end{equation}
    \begin{proof}
        两边都是关系，
        并且
        \begin{align*}
            \opair{x,y}\in Q\SetRestrict(A\cup B)
            &\iff\opair{x,y}\in Q\land(x\in A\lor x\in B) \\
            &\iff(\opair{x,y}\in Q\land x\in A)
            \lor(\opair{x,y}\in Q\land x\in B) \\
            &\iff\opair{x,y}\in(Q\SetRestrict A)\cup(Q\SetRestrict B).
        \end{align*}
        由外延公理得到结论.
    \end{proof}
\end{example}

\begin{example}
    %@see: 《Elements of Set Theory》 P65 Exercise 59.
    设\(A,B,Q\)是集合.
    证明：
    \begin{gather}
        Q \SetRestrict (A \cap B)
        = (Q \SetRestrict A) \cap (Q \SetRestrict B), \\
        Q \SetRestrict (A - B)
        = (Q \SetRestrict A)
        - (Q \SetRestrict B).
    \end{gather}
    \begin{proof}
        对任意\(x,y\)，
        由限制的定义有
        \begin{align*}
            \opair{x,y}\in Q\SetRestrict(A\cap B)
            &\iff\opair{x,y}\in Q\land x\in A\land x\in B \\
            &\iff\opair{x,y}\in(Q\SetRestrict A)\cap(Q\SetRestrict B), \\
            \opair{x,y}\in Q\SetRestrict(A-B)
            &\iff\opair{x,y}\in Q\land x\in A\land x\notin B \\
            &\iff\opair{x,y}\in Q\SetRestrict A
            \land\opair{x,y}\notin Q\SetRestrict B \\
            &\iff\opair{x,y}\in(Q\SetRestrict A)-(Q\SetRestrict B).
        \end{align*}
        各式两边都只含有序对，
        故由外延公理得到所需等式.
    \end{proof}
\end{example}

\begin{example}
    %@see: 《Elements of Set Theory》 P65 Exercise 60.
    设\(A,R,S\)是集合.
    证明：
    \begin{equation}
        (R \circ S) \SetRestrict A = R \circ (S \SetRestrict A).
    \end{equation}
    \begin{proof}
        两边都是关系，
        并且对任意\(x,y\)有
        \begin{align*}
            \opair{x,y}\in(R\circ S)\SetRestrict A
            &\iff x\in A\land(\exists t)[\opair{x,t}\in S\land\opair{t,y}\in R] \\
            &\iff(\exists t)[\opair{x,t}\in S\SetRestrict A\land\opair{t,y}\in R] \\
            &\iff\opair{x,y}\in R\circ(S\SetRestrict A).
        \end{align*}
        由外延公理得到结论.
    \end{proof}
\end{example}

\begin{proposition}\label{theorem:集合论.与逆相等的充分必要条件}
    设\(R\)是关系，
    则\(R^{-1} \subseteq R
    \iff R^{-1} = R
    \iff R \subseteq R^{-1}\).
    \begin{proof}
        先设\(R^{-1}\subseteq R\).
        任取\(\opair{x,y}\in R\)，
        则\(\opair{y,x}\in R^{-1}\subseteq R\)，
        进而\(\opair{x,y}\in R^{-1}\).
        由于\(R\)只含有序对，
        这就证明了\(R\subseteq R^{-1}\)，
        所以\(R=R^{-1}\).

        再设\(R\subseteq R^{-1}\).
        任取\(\opair{x,y}\in R^{-1}\)，
        则\(\opair{y,x}\in R\subseteq R^{-1}\)，
        进而\(\opair{x,y}\in R\).
        因此\(R^{-1}\subseteq R\)，
        仍有\(R=R^{-1}\).

        反过来，
        由\(R=R^{-1}\)立即得到两个方向的包含关系，
        故三个条件等价.
    \end{proof}
\end{proposition}

\begin{example}
    举例说明：在上一个命题中，
    不能把“\(R\)是关系”改成“\(R\)是任意集合”.
    \begin{solution}
        取\(R=\Set{\emptyset}\).
        按有序对的定义，
        \(\opair{x,y}=\Set{\Set{x},\Set{x,y}}\)至少含有元素\(\Set{x}\)，
        所以空集不是有序对.
        因而\(R\)不含任何有序对，
        其逆为\(R^{-1}=\emptyset\).
        于是
        \begin{align*}
            R^{-1} &=\emptyset\subseteq R, \\
            R^{-1} &\neq R, \\
            R &\nsubseteq R^{-1}.
        \end{align*}
        这就给出了原条件过宽时的反例.
    \end{solution}
\end{example}

\begin{example}\label{example:集合论.两个单根集的复合是单根的}
    %@see: 《Elements of Set Theory》 P53 Exercise 17.
    证明：两个单根集的复合还是单根的.
    \begin{proof}
        设\(A,B\)是单根集，
        令\(C=A\circ B\).
        任取\(y\in\ran C\)，
        至少存在一个\(x\)满足\(\opair{x,y}\in C\).
        下面验证这样的\(x\)唯一.

        设\(\opair{x_1,y},\opair{x_2,y}\in C\).
        由复合的定义，
        存在\(t_1,t_2\)满足
        \begin{align*}
            \opair{x_1,t_1}\in B &\land\opair{t_1,y}\in A, \\
            \opair{x_2,t_2}\in B &\land\opair{t_2,y}\in A.
        \end{align*}
        由\(A\)的单根性可得\(t_1=t_2\)，
        再由\(B\)的单根性可得\(x_1=x_2\).
        因此每个\(y\in\ran C\)恰有一个原像，
        故\(C\)单根.
        证明只涉及\(A,B\)中的有序对，
        不需要假设它们是映射.
    \end{proof}
\end{example}

\subsection{有标集族，指标集，子集族}\label{section:集合论.指标集}
%@see: https://math.libretexts.org/Bookshelves/Mathematical_Logic_and_Proof/Book%3A_Mathematical_Reasoning__Writing_and_Proof_(Sundstrom)/05%3A_Set_Theory/5.05%3A_Indexed_Families_of_Sets
\begin{definition}
%@see: 《Elements of Set Theory》 P51
设\(F\)是映射，\(I\)是集合，\(\dom F \supseteq I\).

对于\(\forall i \in I\)，
把\(i\)在映射\(F\)下的像\(F(i)\)记作\(F_i\)，
即\begin{equation*}
	F_i \defeq F(i).
\end{equation*}

把\(F\)在\(I\)上的限制\(F \SetRestrict I\)
称为“一个以\(I\)为指标集的\DefineConcept{有标集族}（an \emph{indexed family of sets} indexed by \(I\)）”，
记作\(\{F_i\}_{i \in I}\)，
即\begin{equation*}
	\{F_i\}_{i\in I}
	\defeq
	F \SetRestrict I.
\end{equation*}

把\(I\)的每一个元素称为一个\DefineConcept{指标}（index）.

把\(I\)称为“有标集族\(\{F_i\}_{i \in I}\)的\DefineConcept{指标集}（indexing set）”.

把\(\ran(F \SetRestrict I)\)称为“有标集族\(\{F_i\}_{i \in I}\)的\DefineConcept{值域}”.
\end{definition}

\begin{example}
    当我们谈到有标集族\(\{R\}_{i \in I}\)时，
    我们指的就是映射\(F = I \times \{R\}\)，
    即对于每一个\(i \in I\)都指定同一个集合\(F(i) = R\).
    \begin{proof}
        按笛卡尔积的定义，
        对任意\(i\in I\)有\(\opair{i,R}\in I\times\Set{R}\).
        若\(\opair{i,y}\in I\times\Set{R}\)，
        则\(y\in\Set{R}\)，
        从而\(y=R\).
        因此\(I\times\Set{R}\)是定义域为\(I\)的映射，
        且\(F(i)=R\)对每个\(i\in I\)成立.
        当\(I=\emptyset\)时，
        该映射就是空映射，
        此时值域为空集而不是\(\Set{R}\).
    \end{proof}
\end{example}

\begin{example}
    设\(\mathscr{A}\)是一个集族，
    那么恒同映射\(\{A\}_{A \in \mathscr{A}} = \Set{ \opair{A,A} \given A \in \mathscr{A} }\)
    就是一个以\(\mathscr{A}\)为指标集的有标集族.
    有时候我们会把\(\{A\}_{A \in \mathscr{A}}\)简记为“有标集族\(\mathscr{A}\)”.
    \begin{proof}
        对每个\(A\in\mathscr{A}\)，
        题给关系中以\(A\)为第一坐标的有序对恰为\(\opair{A,A}\).
        因此它是定义域为\(\mathscr{A}\)、在\(A\)处取值为\(A\)的映射，
        也就是\(I_{\mathscr{A}}\).
        按有标集族的定义，
        以\(\mathscr{A}\)为指标集时，
        其第\(A\)个值恰是集合\(A\)，
        所以它确实表示有标集族\(\{A\}_{A\in\mathscr{A}}\).
    \end{proof}
\end{example}

\begin{definition}
%@see: 《Elements of Set Theory》 P51
%@see: 《点集拓扑讲义（第四版）》（熊金城） P26 定义1.6.1
设有标集族\(\{F_i\}_{i \in I}\).

定义：\begin{equation}
	\bigcup_{i \in I} F_i
	\defeq
	\bigcup\Set{ F_i \given i \in I },
\end{equation}
把它称为“有标集族\(\{F_i\}_{i \in I}\)的并”.

当指标集\(I\)非空时，定义：
\begin{equation}
	\bigcap_{i \in I} F_i
	\defeq
	\bigcap\Set{ F_i \given i \in I },
\end{equation}
把它称为“有标集族\(\{F_i\}_{i \in I}\)的交”.
\end{definition}
应该注意到：
\(\bigcup_{i \in \emptyset} F_i = \emptyset\)，
而\(\bigcap_{i \in \emptyset} F_i\)没有定义.

\begin{theorem}
    %@see: 《点集拓扑讲义（第四版）》（熊金城） P26 定理1.6.1
    设\(\{A_i\}_{i \in I}\)和\(\{B_j\}_{j \in J}\)是两个非空有标集族.
    如果
    \begin{equation*}
        \Set{ A_i \given i \in I }
        = \Set{ B_j \given j \in J },
    \end{equation*}
    则有
    \begin{gather}
        \bigcup_{i \in I} A_i = \bigcup_{j \in J} B_j, \\
        \bigcap_{i \in I} A_i = \bigcap_{j \in J} B_j.
    \end{gather}
    \begin{proof}
        令\(\mathscr{C}=\Set{A_i\given i\in I}=\Set{B_j\given j\in J}\).
        由于\(I,J\)非空，
        \(\mathscr{C}\)也非空.
        对任意\(x\)，
        有
        \begin{align*}
            x\in\bigcup_{i\in I}A_i
            &\iff(\exists C\in\mathscr{C})[x\in C] \\
            &\iff x\in\bigcup_{j\in J}B_j, \\
            x\in\bigcap_{i\in I}A_i
            &\iff(\forall C\in\mathscr{C})[x\in C] \\
            &\iff x\in\bigcap_{j\in J}B_j.
        \end{align*}
        由外延公理得到两个等式.
        因此指标的名称及同一个集合被重复列出的次数，
        都不影响并与交.
    \end{proof}
\end{theorem}

\begin{theorem}
    %@see: 《点集拓扑讲义（第四版）》（熊金城） P27 定理1.6.2
    设\(\{A_i\}_{i \in I}\)是一个非空有标集族，
    \(B\)是一个集合，
    则
    \begin{itemize}
        \item 对于任意\(i_0 \in I\)，有
        \begin{equation*}
            \bigcap_{i \in I} A_i \subseteq A_{i_0} \subseteq \bigcup_{i \in I} A_i;
        \end{equation*}

        \item {\rm\bf 分配律}
        \begin{gather*}
            B \cap \left( \bigcup_{i \in I} A_i \right)
            = \bigcup_{i \in I} \left( B \cap A_i \right), \\
            B \cup \left( \bigcap_{i \in I} A_i \right)
            = \bigcap_{i \in I} \left( B \cup A_i \right);
        \end{gather*}

        \item {\rm\bf 对偶律}
        \begin{gather*}
            B - \left( \bigcup_{i \in I} A_i \right)
            = \bigcap_{i \in I} \left( B - A_i \right), \\
            B - \left( \bigcap_{i \in I} A_i \right)
            = \bigcup_{i \in I} \left( B - A_i \right).
        \end{gather*}
    \end{itemize}
    \begin{proof}
        对固定的\(i_0\in I\)，
        若\(x\in\bigcap_{i\in I}A_i\)，
        则\(x\)属于每个\(A_i\)，
        特别地属于\(A_{i_0}\).
        若\(x\in A_{i_0}\)，
        则\(i_0\)见证\(x\in\bigcup_{i\in I}A_i\).
        所以第一项的两个包含关系成立.

        对任意\(x\)，
        两条分配律分别由下列等价关系得到：
        \begin{align*}
            x\in B\cap\bigcup_{i\in I}A_i
            &\iff x\in B\land(\exists i\in I)[x\in A_i] \\
            &\iff(\exists i\in I)[x\in B\land x\in A_i] \\
            &\iff x\in\bigcup_{i\in I}(B\cap A_i), \\
            x\in B\cup\bigcap_{i\in I}A_i
            &\iff x\in B\lor(\forall i\in I)[x\in A_i] \\
            &\iff(\forall i\in I)[x\in B\lor x\in A_i] \\
            &\iff x\in\bigcap_{i\in I}(B\cup A_i).
        \end{align*}
        两条对偶律分别由量词的否定规则得到：
        \begin{align*}
            x\in B-\bigcup_{i\in I}A_i
            &\iff x\in B\land(\forall i\in I)[x\notin A_i] \\
            &\iff(\forall i\in I)[x\in B\land x\notin A_i] \\
            &\iff x\in\bigcap_{i\in I}(B-A_i), \\
            x\in B-\bigcap_{i\in I}A_i
            &\iff x\in B\land(\exists i\in I)[x\notin A_i] \\
            &\iff(\exists i\in I)[x\in B\land x\notin A_i] \\
            &\iff x\in\bigcup_{i\in I}(B-A_i).
        \end{align*}
        其中第一条对偶律把\(x\in B\)移入全称量词时，
        使用了\(I\neq\emptyset\).
        最后分别应用外延公理即可.
    \end{proof}
\end{theorem}

\begin{definition}
%@see: 《点集拓扑讲义（第四版）》（熊金城） P28
设\(X\)是一个集合，
\(\{A_i\}_{i \in I}\)是一个有标集族，
且满足\begin{equation*}
	(\forall i \in I)
	[A_i \subseteq X],
\end{equation*}
则称“\(\{A_i\}_{i \in I}\)是集合\(X\)的一个\DefineConcept{子集族}”.
\end{definition}

\begin{theorem}
    %@see: 《点集拓扑讲义（第四版）》（熊金城） P28 定理1.6.3
    设\(R\)是集合\(X\)与集合\(Y\)之间的一个关系，
    则对于集合\(X\)的任何一个非空子集族\(\{A_i\}_{i \in I}\)，
    有
    \begin{gather*}
        R\ImageOfSetUnderRelation{\bigcup_{i \in I} A_i}
        = \bigcup_{i \in I} R\ImageOfSetUnderRelation{A_i}, \\
        R\ImageOfSetUnderRelation{\bigcap_{i \in I} A_i}
        \subseteq \bigcap_{i \in I} R\ImageOfSetUnderRelation{A_i}.
    \end{gather*}
    \begin{proof}
        任取\(y\in Y\)，
        由关系的像的定义可得
        \begin{align*}
            y\in R\ImageOfSetUnderRelation{\bigcup_{i\in I}A_i}
            &\iff(\exists x\in\bigcup_{i\in I}A_i)[\opair{x,y}\in R] \\
            &\iff(\exists i\in I)(\exists x\in A_i)[\opair{x,y}\in R] \\
            &\iff y\in\bigcup_{i\in I}R\ImageOfSetUnderRelation{A_i}.
        \end{align*}
        两边均为\(Y\)的子集，
        因此并的等式成立.

        若\(y\in R\ImageOfSetUnderRelation{\bigcap_{i\in I}A_i}\)，
        取\(x\in\bigcap_{i\in I}A_i\)满足\(\opair{x,y}\in R\).
        对每个\(i\in I\)，
        这个\(x\)都属于\(A_i\)，
        从而\(y\in R\ImageOfSetUnderRelation{A_i}\).
        因此\(y\in\bigcap_{i\in I}R\ImageOfSetUnderRelation{A_i}\)，
        得到交的包含关系.
    \end{proof}
\end{theorem}

\begin{theorem}
    %@see: 《点集拓扑讲义（第四版）》（熊金城） P28 定理1.6.4
    设\(X\)和\(Y\)是两个集合，映射\(f\colon X \to Y\)，
    则对于集合\(Y\)的任何一个非空子集族\(\{B_j\}_{j \in J}\)，
    有
    \begin{gather*}
        f^{-1}\ImageOfSetUnderRelation{\bigcup_{j \in J} B_j}
        = \bigcup_{j \in J} f^{-1}\ImageOfSetUnderRelation{B_j}, \\
        f^{-1}\ImageOfSetUnderRelation{\bigcap_{j \in J} B_j}
        = \bigcap_{j \in J} f^{-1}\ImageOfSetUnderRelation{B_j}.
    \end{gather*}
    \begin{proof}
        下列等式两边均为\(X\)的子集，
        所以只需检验任意\(x\in X\)的归属.
        由原像的定义可得
        \begin{align*}
            x\in f^{-1}\ImageOfSetUnderRelation{\bigcup_{j\in J}B_j}
            &\iff f(x)\in\bigcup_{j\in J}B_j \\
            &\iff(\exists j\in J)[f(x)\in B_j] \\
            &\iff x\in\bigcup_{j\in J}f^{-1}\ImageOfSetUnderRelation{B_j}, \\
            x\in f^{-1}\ImageOfSetUnderRelation{\bigcap_{j\in J}B_j}
            &\iff f(x)\in\bigcap_{j\in J}B_j \\
            &\iff(\forall j\in J)[f(x)\in B_j] \\
            &\iff x\in\bigcap_{j\in J}f^{-1}\ImageOfSetUnderRelation{B_j}.
        \end{align*}
        由外延公理得到两个等式.
        这里\(f^{-1}\ImageOfSetUnderRelation{B_j}\)是集合的原像，
        不是对\(B_j\)作逆映射求值，
        所以无需假设\(f\)可逆.
    \end{proof}
\end{theorem}

\subsection{映射空间}
对于任意给定的集合\(A,X\)，定义：\begin{equation*}
	X^A \defeq \Set{ F \given F\ \text{是从\(A\)到\(X\)的映射} }.
\end{equation*}
我们把\(X^A\)称为“从\(A\)到\(X\)的\DefineConcept{映射空间}（function space）”.

因为\(F\colon A \to X\)必有\(F \subseteq A \times X\)，\(F \in \Powerset(A \times X)\)，
所以我们可以对集合\(\Powerset(A \times X)\)利用子集公理，构造包括全部从\(A\)到\(X\)的映射的集合.

%之所以采取这种表记方式，
%是因为当\(A\)和\(X\)是有限集，且\(\abs{A}=a,\abs{X}=x\)时，
%\(\abs{X^A}=x^a\).

容易看出，对于非空集合\(A\)，总有\(\emptyset^A = \emptyset\)；
这是因为没有哪个映射会同时有非空的定义域和空的值域.
另一方面，对于任意集合\(A\)，总有\(A^\emptyset = \Set{\emptyset}\)；
这是因为“空映射”\(\emptyset\colon \emptyset \to A\)的存在，
空映射是唯一的以空集为定义域的映射.
作为特例，我们还有\(\emptyset^\emptyset=\Set{\emptyset}\).

\subsection{映射的四则运算}
\begin{definition}
设\(D\)是集合，
\(\mathbb{K}\)为\(\mathbb{R}\)或\(\mathbb{C}\)，
\(F,G\colon D\to\mathbb{K}\)，
\(k\in\mathbb{K}\).
\begin{itemize}
	\item 把\begin{equation*}
		\Set{ \opair{u,kv} \given \opair{u,v} \in F }
	\end{equation*}称为“\(F\)的\(k\)~\DefineConcept{倍}”，
	记作\(k F\).
	\item 把\begin{equation*}
		\Set{ \opair{u,v+w} \given \opair{u,v} \in F, \opair{u,w} \in G }
	\end{equation*}称为“\(F\)和\(G\)的\DefineConcept{和}”，
	记作\(F + G\).
	\item 把\begin{equation*}
		\Set{ \opair{u,vw} \given \opair{u,v} \in F, \opair{u,w} \in G }
	\end{equation*}称为“\(F\)和\(G\)的\DefineConcept{积}”，
	记作\(F G\).
	\item 把\begin{equation*}
		\Set{ \opair{u,v/w} \given \opair{u,v} \in F, \opair{u,w} \in G,\ w\neq0 }
	\end{equation*}为“\(F\)和\(G\)的\DefineConcept{商}”，
	记作\(F/G\)或\(\frac{F}{G}\).
\end{itemize}
商映射的定义域是\(\Set{u\in D\given G(u)\neq0}\).
只有当\(G\)在\(D\)上恒不为零时，
商映射才与\(F,G\)具有同一个定义域.
\end{definition}
